%%%%%%%%%%  scan idx 97  |  folio 86  %%%%%%%%%%

\underline{Teorema 3.12.}

\begin{enumerate}[label=(\arabic*)]
\item Si
\[
 0\longrightarrow M'\longrightarrow M\longrightarrow M''
 \longrightarrow0
\]
es una sucesión exacta de $A$-módulos, entonces la sucesión
\begin{multline*}
0\longrightarrow D^0(A|R,M')\longrightarrow D^0(A|R,M)
\longrightarrow D^0(A|R,M'')\longrightarrow D^1(A|R,M')\\
\longrightarrow D^1(A|R,M)\longrightarrow D^1(A|R,M'')
\longrightarrow\cdots\longrightarrow D^n(A|R,M')\\
\longrightarrow D^n(A|R,M)\longrightarrow D^n(A|R,M'')
\longrightarrow\cdots
\end{multline*}
es exacta.

\item Sean $S\longrightarrow R\longrightarrow A$ homomorfismos de
anillos. Entonces, para todo $A$-módulo $M$, la sucesión
\begin{multline*}
0\longrightarrow D^0(A|R,M)\longrightarrow D^0(A|S,M)
\longrightarrow D^0(R|S,M)\longrightarrow D^1(A|R,M)\\
\longrightarrow D^1(A|S,M)\longrightarrow D^1(R|S,M)
\longrightarrow\cdots\longrightarrow D^n(A|R,M)\\
\longrightarrow D^n(A|S,M)\longrightarrow D^n(R|S,M)
\longrightarrow\cdots
\end{multline*}
es exacta.

\item Sean $A$ y $S$ $R$-álgebras tales que
$\operatorname{Tor}_i^R(A,S)=0$ para todo $i>0$. Entonces, para todo
$A\otimes_RS$-módulo $M$, el morfismo canónico
\[
 D^i(A\otimes_RS|S,M)\longrightarrow D^i(A|R,M)
\]
es un isomorfismo para todo $i$. En particular, si $A$ es plano sobre
$R$, entonces
\[
 D^i(A\otimes_RS|S,M)\xrightarrow{\ \sim\ }D^i(A|R,M)
\]
para todo $i$, toda $R$-álgebra $S$ y todo $A\otimes_RS$-módulo $M$.

\item Sea $S$ un subconjunto multiplicativo de $A$. Para todo
$S^{-1}A$-módulo $N$, el morfismo canónico
\[
 D^i(S^{-1}A|R,N)\longrightarrow D^i(A|R,N)
\]
es un isomorfismo para todo $i$. Por consiguiente, para todo $A$-módulo
$M$, hay un isomorfismo canónico
\[
 D^i(S^{-1}A|R,S^{-1}M)
 \xrightarrow{\ \sim\ }S^{-1}D^i(A|R,M)
\]
para todo $i$.

\item Sean $A$ y $B$ $R$-álgebras tales que
$\operatorname{Tor}_i^R(A,B)=0$ para $0<i\leq n$. Entonces, para todo
$A\otimes_RB$-módulo $M$, el morfismo canónico
\[
 D^i(A\otimes_RB|R,M)
 \longrightarrow D^i(A|R,M)\oplus D^i(B|R,M)
\]
es un isomorfismo para $i\leq n$. En particular, supongamos que
$\operatorname{Tor}_i^R(A,A)=0$ para todo $i>0$ y que la diagonal
\[
 \Delta:\operatorname{Spec}(A)\longrightarrow
 \operatorname{Spec}(A)\times_{\operatorname{Spec}(R)}
 \operatorname{Spec}(A)
\]
es una inmersión abierta; por ejemplo, supongamos que $A$ es
\emph{étale} sobre $R$. Entonces
\[
 D^i(A|R,M)=0
\]
para todo $i>0$ y todo $A$-módulo $M$.

%%%%%%%%%%  scan idx 98  |  folio 87  %%%%%%%%%%

\item Sea $R$ noetheriano, sea $A$ una $R$-álgebra esencialmente de
tipo finito y sea $E$ un $A$-módulo plano. Entonces el morfismo canónico
\[
 E\otimes_AD^i(A|R,M)
 \xrightarrow{\ \sim\ }
 D^i(A|R,E\otimes_AM)
\]
es un isomorfismo para todo $i\geq0$ y todo $A$-módulo $M$.

\item Si $R$ es noetheriano, $A$ es una $R$-álgebra esencialmente de
tipo finito y $M$ es un $A$-módulo de tipo finito, entonces
$D^i(A|R,M)$ es un $A$-módulo de tipo finito para todo $i$.

\item Sea $P$ un álgebra de polinomios sobre $R$, y sea exacta
\[
 0\longrightarrow I\longrightarrow P\longrightarrow A
 \longrightarrow0.
\]
Entonces
\[
 D^0(A|R,M)=\operatorname{Der}_R(A,M),
\]
\begin{align*}
D^1(A|R,M)
 &=\operatorname{Coker}\bigl(
   \operatorname{Der}_R(P,M)\\[-.25ex]
 &\hspace{5em}\longrightarrow
   \operatorname{Hom}_A(I/I^2,M)\sisi{}{\bigr)}\\
 &=\left\{
   \begin{array}{c}
   \text{clases de isomorfismo de extensiones}\\
   \text{conmutativas de }R\text{-álgebras de }A\text{ por }M
   \end{array}\right\},\\
D^2(A|R,M)
 &=\operatorname{Coker}\bigl(
   H^1(K_I,M)\longrightarrow
   \operatorname{Hom}_A(H_1(K_I),M)\bigr).
\end{align*}
\end{enumerate}
